\begin{problem}{}{}{Real Mathematical Analysis, Pugh, Ex.2 p.198}
A function $f : (a,b) \to \mathbb{R}$ satisfies a \textbf{H\"older condition}
of order $\boldsymbol{\alpha}$ if $\alpha > 0$, and for some constant $H$
and all $u,x \in (a,b)$ we have
\[
    |f(u) - f(x)| \le H|u - x|^\alpha
\]
The function is said to be $\alpha$-H\"older, with $\alpha$-H\"older
constant $H$. (The terms ``Lipschitz function of order $\alpha$'' and
``$\alpha$-Lipschitz function'' are sometimes used with the same meaning.)

\begin{enumerate}
    \item[(a)] Prove that an $\alpha$-H\"older function defined on $(a,b)$
          is uniformly continuous and infer that it extends uniquely to a
          continuous function defined on $[a, b]$. Is the extended function
          $\alpha$-H\"older?

    \item[(b)] What does $\alpha$-H\"older continuity mean when $\alpha = 1$?

    \item[(c)] Prove that $\alpha$-H\"older continuity when $\alpha > 1$
          implies that $f$ is constant.
\end{enumerate}
\end{problem}

\begin{solution}
    \begin{itemize}
        \item[(a)]
              Let $\varepsilon > 0$. Choose $\delta = \sqrt[\alpha]{\frac{\varepsilon}{H}}$. Now, for every $u, x \in (a, b)$ and $|u - x| < \delta$ implies
              \[
                  |f(u) - f(x)| \leq H|u - x|^\alpha < H\delta^\alpha = H \frac{\varepsilon}{H} = \varepsilon.
              \]
              Thus, $f$ is uniformly continuous on $(a, b)$.

              \begin{remark}{}
                  Uniform continuity implies that if $(x_n)$ is Cauchy then $(f(x_n))$ is also a Cauchy sequence.
              \end{remark}

              Let $(x_n)$ be a sequence in $(a, b)$, $(x_n) \to a$. It follows that $(f(x_n))$ also converges
              to a finite value $L$ since $\mathbb{R}$ is complete. It makes sense to define $\lim_{x \to a^+} f(x) = L$
              by Theorem 4.2.3 (\textit{Understanding Analysis}, Abbott).

              For $g$ to be an extension of $f$ that is continuous on $[a, b]$. We must define
              $g(a) = \lim_{x \to a^+} f(x) = L$. Likewise for endpoint $b$. Since this is the
              only way to construct such a $g$, $g$ is unique.

              To check if $g$ is $\alpha$-Hölder, we need to check:
              \begin{itemize}
                  \item[(1):] $|g(a) - g(x)| = |L - f(x)| \leq H|a - x|^\alpha$, for all $x \in (a, b)$

                        This is true since for any $t, x \in (a, b)$, we have:
                        \begin{alignat*}{2}
                             &                & |f(t) - f(x)|                  & \leq H|t - x|^\alpha                  \\
                             & \implies \quad & \lim_{t \to a^+} |f(t) - f(x)| & \leq \lim_{t \to a^+} H|t - x|^\alpha \\
                             & \implies \quad & |L - f(x)|                     & \leq H|a - x|^\alpha
                        \end{alignat*}

                  \item[(2):] $|g(b) - g(x)| \leq H|b - x|^\alpha$ \quad (true for similar reason)
                  \item[(3):] $|g(b) - g(a)| \leq H|b - a|^\alpha$ \quad (take limit of (2) as $x \to a^+$)
              \end{itemize}
        \item[(b)]
              If $\alpha = 1$, the $\alpha$-Hölder condition becomes:
              $$ |f(u) - f(x)| \le H|u - x| $$

              Geometrically, this means the absolute value of the slope
              of the secant line between any two points is bounded by $H$.
        \item[(c)]
              Let $\alpha > 1$. Let $u,x \in (a,b)$, $u \neq x$. We have:
              \begin{alignat*}{2}
                   &                & |f(u) - f(x)|                                           & \leq H|u - x|^\alpha                      \\
                   & \implies \quad & \left| \frac{f(u) - f(x)}{u - x} \right|                & \leq H|u - x|^{\alpha - 1}                \\
                   & \implies \quad & \lim_{x \to u} \left| \frac{f(u) - f(x)}{u - x} \right| & \leq \lim_{x \to u} H|u - x|^{\alpha - 1} \\
                   & \implies \quad & |f'(u)|                                                 & \leq 0                                    \\
                   & \implies \quad & f'(u)                                                   & = 0, \quad \forall u \in (a,b)
              \end{alignat*}
              This implies $f$ is constant on $(a,b)$.
    \end{itemize}
\end{solution}